\subsection{Local-linear-span and recurrence methods}

The budgeted family $\mathsf{RowRec}$ sharpened in
\cref{sec:resource-vector} requires that new numerical directions are
generated from previously stored directions by charged local
row actions, supported matrix--vector products, diagonal or fixed splitting
operations, and scalar linear combinations.  The definition must specify:

\begin{enumerate}[label=(\alph*)]
\item whether a round may apply all exposed rows or only selected rows;
\item whether arbitrary diagonal scaling is allowed;
\item whether momentum, conjugate directions, restarts, and reorthogonalization
      are allowed;
\item how many persistent scalars may be stored per vertex or incidence;
\item whether the support of a direction may be masked or projected;
\item whether a graph-dependent preconditioner may be built and, if so, how
      its construction and application are charged.
\end{enumerate}

Without item (f), the class is not stable: a method described syntactically as
preconditioned CG could encode the entire restricted inverse in its
preconditioner and have the power of a response solver.  The formal collapse
is \cref{lem:free-preconditioner-collapse}.

The classical square-root condition-number barrier is relevant to a suitable
linear-span class.  It is not automatically a lower bound for
$\mathsf{Adj}$.
